\scanpage{39}
% archive sheet number: 38
connaissance de $\underline{D}^{*}(G_{\unsure{o|0}})$ permet de réconstituer $G_o$. Ainsi, \deleted{\unsure{a}}dans le cas général, la connaissance de $\underline{D}^{*}(G_o)=M$ comme F-V-cristal sur $S_o$ permet de retrouver les fibres de $G_o$ en des corps parfaits au dessus de $S_o$ ; cela justifie donc le sentiment que pour l'essentiel, on a saisi avec le cristal $M=\underline{D}^{*}(G_o)$ la \deleted{\unsure{f}}famille des $G_{o\bar{s}}$. Une question importante \inserted{\handwritten{en ce sens}} que je n'ai pas résolue est :

\uline{Problème} 1 : Le foncteur de Dieudonné sur une base $S_o$ de car. $p$ est-il pleinement fidèle ? \deleted{Quel est}\deleted{\inserted{\unsure{le|la}}} \deleted{\unsure{s}}Son image essentielle \deleted{\unsure{e|S}}est-elle formée des F-V-cristaux ``admissibles'' définis plus bas ?

Pour résoudre cette question, il me semble qu'il faudrait arriver à formuler une théorie de Dieudonné s\deleted{\unsure{r|u}}ur \deleted{des bases de car $p$} $S_o$ pour des schémas en \inserted{$p$-}groupes finis localement libres sur $S_o$ quelconques, ou du moins pour ceux plats sur un $\handwritten{\Lambda}_n$, en trouvant une équivalence entre \deleted{cette} la catégorie de\deleted{s} ces groupes sur $S_o$, et une catégorie de F-V-cristaux sur $S_o$ \uline{muni de structures supplémentaires}\inserted{\handwritten{$^{(1)}$}}\margin{\handwritten{sans doute inutile si $S_o$ est parfait}} que je n'ai pas réussi à dégager encore, et qui très probablement impliqueront la donnée d'une \uline{filtration} sur un objet \inserted{convenable} d'une catégorie dérivée \ldots
